\documentclass[11pt,a4paper]{article}
\usepackage[margin=1.1in]{geometry}
\usepackage{amsmath,amssymb}
\usepackage{mathtools}
\usepackage{booktabs}
\usepackage{algorithm,algorithmicx,algpseudocode}
\usepackage{enumitem}
\usepackage{mdframed}
\usepackage[colorlinks=true,linkcolor=blue]{hyperref}

\newcommand{\R}{\mathbb{R}}
\newcommand{\poly}{\varphi}

\title{\bfseries Newton--Schulz Coefficients\\
for the Muon Optimizer}
\author{}
\date{}

\begin{document}
\maketitle

\section{Background}

\subsection{Newton--Schulz Orthogonalization}

The Muon optimizer~\cite{jordan2024} uses Newton--Schulz (NS) iteration to
orthogonalize gradient momentum matrices.  For a weight matrix
$W \in \R^{m \times n}$, the output of SGD-momentum, denoted
$G_t \in \R^{m \times n}$, is fed into the iteration as a post-processing
step:
\[
X_0 = \frac{G_t}{\|G_t\|_F}, \qquad
X_{k+1} = a X_k + b X_k X_k^\top X_k + c (X_k X_k^\top)^2 X_k,
\]
and the weights are updated as $W_{t+1} = W_t - \eta X_5$.

The iteration acts on the singular values of $G_t$ through the scalar
polynomial
\begin{equation}\label{eq:poly}
\poly(x) = a x + b x^3 + c x^5,
\end{equation}
so that if $G_t / \|G_t\|_F = P \Sigma Q^\top$, then the NS output after
$k$ steps is $P \, \poly^{(k)}(\Sigma) \, Q^\top$, where
$\poly^{(k)}$ denotes $k$-fold composition.

\subsection{Existing Coefficient Sets}

Two coefficient sets are provided as baselines:

\begin{center}
\begin{tabular}{lrrrl}
\toprule
& $a$ & $b$ & $c$ & Source \\
\midrule
Standard NS & 1.500 & $-0.500$ & 0.000 & Classical degree-3 iteration \\
Jordan NS   & 3.445 & $-4.775$ & 2.032 & Numerical search~\cite{jordan2024} \\
\bottomrule
\end{tabular}
\end{center}

Standard NS satisfies the exact constraints $\poly(1)=1$ and
$\poly'(1)=0$, guaranteeing quadratic convergence near 1.  However,
$a = \poly'(0) = 1.5$ provides weak pull from zero: after 5 iterations,
a singular value of $10^{-3}$ reaches only $8 \times 10^{-3}$.

Jordan NS relaxes the fixed-point constraint ($\poly(1) \approx 0.70$) in
exchange for a larger $a = 3.44$.  The same initial value reaches $0.47$ in 5 iterations, yielding substantially faster convergence in practice,
at the cost of increased step-to-step loss variance.

\subsection{The Problem}

Fixed coefficients cannot be optimal throughout training.  In early
stages, gradients are high-rank and noisy, favoring conservative $a$.
In late training, gradients become low-rank, and aggressive $a$
accelerates convergence.  Moreover, $a=3.44$ was found for a specific
setting, and there is no guarantee it is optimal for other problems.

This exercise addresses two questions:
\begin{enumerate}[nosep]
    \item What is the maximal stable value of $a$, and what $(b,c)$
          support it?
    \item Can $a$ be self-tuned during training via a stability
          feedback signal, avoiding manual coefficient selection?
\end{enumerate}

\section{Stability Constraints}

A candidate $(a,b,c)$ triple must satisfy two conditions.

\subsection{5-Iteration Band}

After exactly $K=5$ NS iterations, every initial singular value
$x \in [\varepsilon, 1]$ where $\varepsilon = 10^{-3}$ must land in
$[0.5, 1.5]$ (this range can be modified):
\begin{equation}\label{eq:band}
\forall x \in [\varepsilon, 1]: \quad 0.5 \leq \poly^{(5)}(x) \leq 1.5.
\end{equation}
The lower bound ensures effective orthogonalization; the upper bound
prevents numerical overflow.

Standard NS fails this constraint:
$\poly^{(5)}(\varepsilon) \approx 0.008 < 0.5$.
Jordan NS is on the borderline: $\poly^{(5)}_{\min} \approx 0.47$. 

\subsection{Orbit Boundedness}

No intermediate iterate may escape (this range can be modified too):
\begin{equation}\label{eq:orbit}
\forall x \in [\varepsilon, 1],\; k = 1,\dots,5:
\quad |\poly^{(k)}(x)| \leq 2.
\end{equation}
If $\poly(x) > 2$, subsequent iterates may exceed the floating-point range,
producing \texttt{NaN} quickly.  This constraint eliminates coefficient sets (such as
$a \approx 4.5$) that would otherwise appear admissible from the final-iterate band alone.

\section{Tasks}

\subsection*{Task 1: Stability Check}

\textbf{File:} \texttt{muon/newton\_schulz.py}\\[2pt]
\textbf{Function:} \texttt{is\_stable(a, b, c, epsilon=1e-3, n\_pts=200)}

Given $(a,b,c)$, verify constraints~\eqref{eq:band}--\eqref{eq:orbit}.
Sample $n\_pts$ log-spaced points in $[\varepsilon, 1]$, iterate the
polynomial 5 times for each, and return \texttt{True} iff all points pass.
\begin{mdframed}
    \texttt{pytest tests/test\_coeff\_search.py::TestIsStable -v}
\end{mdframed}

\subsection*{Task 2: Random Coefficient Search}

\textbf{File:} \texttt{muon/coeff\_search.py}\\[2pt]
\textbf{Function:} \texttt{search\_coefficients(n\_samples=50000, ...)}

Randomly sample $(a,b,c)$ from specified ranges, evaluate each with
\texttt{is\_stable}, and return the triple with the largest $a$.

\begin{mdframed}
    \texttt{pytest tests/test\_coeff\_search.py::TestCoeffSearch -v}
\end{mdframed}

\subsection*{Task 3: Maximal Stable Coefficient}

\textbf{File:} \texttt{muon/coeff\_search.py}\\[2pt]
\textbf{Function:} \texttt{find\_fast\_coeffs(n\_samples=200000)}

Find the stable $(a,b,c)$ with the maximum $a$.  The result should exceed Jordan's $a=3.44$.

\begin{mdframed}
\texttt{pytest tests/test\_coeff\_search.py -k find\_fast -v}
\end{mdframed}

\subsection*{Task 4: Coefficient Lookup Table}

\textbf{File:} \texttt{muon/coeff\_search.py}\\[2pt]
\textbf{Function:} \texttt{build\_coeff\_table(a\_values)}

For each target $a$ in the input list, run a focused search to find the best stable $(b,c)$ near that $a$.  Return a dictionary
$\{a_{\rm target} \mapsto (a,b,c)\}$.

The adaptive optimizer's default table contains only two entries
($a=1.5$ and $a=3.44$).  When the adaptive scheduler requests $a=4.0$,
nearest-key lookup returns $a=3.44$, and the optimizer never uses
aggressive coefficients.  This task populates the table with intermediate entries so the adaptive method can function correctly.

\subsection*{Task 5: Adaptive Coefficient Scheduling}

\textbf{File:} \texttt{muon/adaptive\_optimizer.py}\\[2pt]
\textbf{Method:} \texttt{AdaptiveMuon.step(loss\_val)}

Implement loss-stability feedback:

\begin{algorithm}[H]
\caption{Adaptive coefficient scheduling}
\begin{algorithmic}[1]
\State $a \leftarrow 1.5$; \quad $\overline{\ell}_t$ $\leftarrow$ None; \quad stable\_ct $\leftarrow 0$
\For{each step with loss $\ell_t$}
    \If{$\overline{\ell}_t$ is None} \State $\overline{\ell}_t$ $\leftarrow \ell_t$
    \Else
        \State $\overline{\ell}_t$ $\leftarrow 0.99 \cdot {}$$\overline{\ell}_t$ ${}+ 0.01 \cdot \ell_t$
        \If{$\ell_t > 3.0 \cdot {}$$\overline{\ell}_t$}\qquad \text{//\textit{NS iteration too aggressive}}
            \State $a \leftarrow \max(a - 0.2,\; 1.5)$
            \State stable\_ct $\leftarrow 0$
        \Else
            \State stable\_ct $\leftarrow {}$stable\_ct ${}+ 1$
            \If{stable\_ct $> 100$}\qquad \text{//\textit{NS iteration too conservative}}
                \State $a \leftarrow \min(a + 0.05,\; 4.0)$
            \EndIf
        \EndIf
    \EndIf
    \State Look up $(a,b,c)$ from table; apply NS; update weights
\EndFor
\end{algorithmic}
\end{algorithm}

\begin{mdframed}
    \texttt{pytest tests/test\_adaptive.py -v}
\end{mdframed}

\subsection*{Task 6: Experiment}

Run \texttt{python run.py}.  This trains an MLP on random oscillatory
functions and compares standard NS, Jordan NS, the fastest coefficients
found by your search, and your adaptive optimizer.  Output:
\texttt{adaptive\_ns\_comparison.png}.

\section{Open Discussion}

\begin{enumerate}
    \item Standard NS satisfies $\poly(1)=1$ and $\poly'(1)=0$ exactly, yet
          $\poly^{(5)}(\varepsilon) \approx 0.008$.  What does this reveal
          about the tension between asymptotic optimality near $x=1$ and
          practical convergence for $x \ll 1$?

    \item Jordan's coefficients violate $\poly(1)=1$ by about $0.30$.
          How far can this constraint be relaxed before the iteration
          becomes unstable?  What mechanism couples $a$ to the allowable deviation
          in $\poly(1)$?

    \item The adaptive method starts at $a=1.5$ and climbs toward $a=4.0$.
          Would initializing directly at $a=4.0$ work?  If not, why?

    \item The spike metric counts epochs where loss exceeds $3\times$
          the running minimum.  For aggressive coefficients, this count
          can reach several hundred.  Are these instabilities
          or artifacts of aggressive exploration?  Propose a metric
          that distinguishes the two.

    \item Could $a$ be adapted per-layer rather than globally?  What signal would inform per-layer adaptation?
\end{enumerate}

\begin{thebibliography}{9}
\bibitem{jordan2024} K.~Jordan.  Muon: An optimizer for hidden layers.
  \texttt{kellerjordan.github.io/posts/muon/}, 2024.
\bibitem{amsel2025} N.~Amsel, T.~Persson, C.~Musco, R.~Gower.
  The Polar Express: optimal matrix sign methods.  \textit{arXiv:2505.16932}, 2025.
\bibitem{liu2025} J.~Liu, J.~Su, et al.  Muon is scalable for LLM training.
  \textit{arXiv:2502.16982}, 2025.
\end{thebibliography}

\end{document}
